\section{问题分析}

本题研究系泊系统的设计。通过确定锚链的型号、长度和重物球质量，使浮标的吃水深度、游动区域及钢桶倾斜角尽可能小，从而保证水声通信设备具有较好的工作效果。

\subsection{问题一分析}

问题一选用 II 型电焊锚链 (22.05\,\mathrm{m})，重物球质量为 (1200\,\mathrm{kg})。传输节点布放在水深 (18\,\mathrm{m})、海床平坦、海水密度为 (1.025\times10^3\,\mathrm{kg/m^3}) 的海域；在海水静止条件下，分别计算海面风速为 (12\,\mathrm{m/s}) 和 (24\,\mathrm{m/s}) 时钢桶和各节钢管的倾斜角度、锚链形状、浮标吃水深度和游动区域。

针对该问题，依次对浮标、钢管、钢桶和重物球、锚链进行受力分析。在问题一的理想条件下，上述物体满足受力平衡，并据此建立方程。对于钢管、钢桶和锚链等刚性物体，进一步建立力矩平衡方程，求解每节钢管、钢桶和各段电焊锚链与海平面的夹角。各段物体在竖直方向投影的总和为海水深度；当钢管、钢桶和各段电焊锚链的竖直投影长度之和满足给定水深时，相应的夹角和吃水深度即为所求。利用钢管、钢桶及各段锚链的倾斜角度，可进一步确定电焊锚链的形状和浮标的游动范围。

\subsection{问题二分析}

对于问题二，根据题目所给数据信息及要求，可利用问题一中的数学模型，求解海面风速为 (36\,\mathrm{m/s}) 时钢桶和各节钢管的倾斜角度、链条形状和浮标游动区域。需要通过调节重物球质量，使钢桶倾斜角不超过 (5^\circ)，最后一节锚链与海床的夹角不超过 (16^\circ)。因此，以钢桶倾斜角和最后一节锚链与海床夹角为约束条件，逐步增加重物球质量，采用二分法对海水深度进行求解；当计算值与实际深度相等时，对应的重物球质量即为所需的最小质量。
